\section{Limitations and Threats to Validity}

\subsection{Threats to Internal Validity}

The representation-quality trade-off currently rests on a fixed-protocol linear
probe. Additional probes, such as kNN or a small MLP, should test whether the
ordering is probe-dependent. The curvature measurements use power iteration on
fixed batches, so sensitivity to batch identity, batch size, and iteration count
should be reported. The $1/2$-saturation diagnostic is demonstrated at one site;
claims about generality should remain cautious until measured across multiple
sites.

The stop-gradient intervention changes the gradient path from CE to $W_1$ while
holding data, labels, architecture, and CE training fixed. This makes it a clean
causal knob, but the argument should be stated explicitly.

\subsection{Limitations of Scope}

The experiments use MNIST and small networks. The paper makes no SOTA claims and
does not offer a general theory of deep-network conditioning. The formal theory
gives sufficient conditions, not a universal if-and-only-if. LayerNorm is
included empirically but omitted from the clean proposition unless folded into
the constants. No convergence proof is provided for the composed system.
